% ============================================================
% Laborbuch – D0rag3on (Noel Neuhoff)
% SafetySpot Projektdokumentation
% Formatierung nach IIRu5hII-Vorlage, Styling per safetyspot.sty
% ============================================================

\ssSetDocTitle{Laborbuch}
\ssSetKunde{Bildungsministerium NRW}
\ssSetProjekt{SafetySpot}
\ssSetAutoren{D0rag3on (Noel Neuhoff)}
\ssSetStand{29.06.2026}
\maketitle
\setcounter{section}{0}

% ============================================================
\section{Grundlage und Projektkontext}
% ============================================================

SafetySpot ist eine mobile Lern-App für Schulen, bei der Schüler interaktiv
Gefahrensituationen aus dem Schulalltag bearbeiten und Punkte sammeln.
Das Gesamtprojekt gliedert sich in drei Module:
\texttt{ssbackend} (Spring-Boot-REST-API), \texttt{ssmobile}
(Android-Kotlin/Compose-App) und \texttt{sscommon} (geteilte Bibliothek).

Mein persönlicher Schwerpunkt lag auf dem Backend:
Aufbau der Projektstruktur, JWT-Authentifizierung, globale Fehlerbehandlung,
School-/Class-Management sowie Umgebungsvariablen-Handling.
Alle Arbeiten wurden im Repository \texttt{safety-spot/SafetySpot} auf GitHub
durchgeführt.

\textit{Hinweis:} Die frühen Commits (April bis Mitte Juni 2026) enthielten
teilweise mehrere Arbeitsschritte in einem einzigen Commit.
Wo nötig, wurden die Einträge sachlich aufgeteilt und mit Querverweis auf den
jeweiligen Git-Commit versehen.

% ============================================================
\section{Ausgewertete Unterlagen / Anlagen}
% ============================================================

\begin{ssunterlagetabelle}
Projektvision (docs/v0) &
  Ziele, Ausgangssituation, Stakeholder, Qualitätsanforderungen und Rahmenbedingungen. \\
Informationsmodell (docs/v0) &
  Persistente Daten: Benutzer, Klassen, Szenarien, Fortschritt, Punkte, Leaderboards
  und Admin-/Lizenzdaten. \\
Klassendiagramm (docs/v1) &
  Fachliche Klassen- und Beziehungsstruktur. \\
Entscheidungslog (docs/v0) &
  Begründungen für Gradle, Spring Boot, MariaDB und Azure. \\
spec/issues/ISSUE-001–011 &
  Geplante Backend-Aufgaben von Foundation über Auth bis Admin-Shell. \\
spec/plan.md &
  Detaillierter Backend-Implementierungsplan (Pakete, APIs, Auth-Strategie). \\
GitHub-Commit-Historie &
  Eigene Commits und Merge-Requests (PRs \#29, \#30, \#31, \#45, \#55). \\
\end{ssunterlagetabelle}

% ============================================================
\section{Technischer Arbeitsstand}
% ============================================================

\begin{ssstandtabelle}
Build-System &
  Migration von Maven auf Gradle (Spring-Boot-Plugin); einheitlich mit ssmobile und sscommon. \\
Paketstruktur &
  Basis-Package \texttt{spot.safety.ssbackend}; Pakete: \texttt{attempt}, \texttt{auth},
  \texttt{dto}, \texttt{enums}, \texttt{exception}, \texttt{leaderboard}, \texttt{progress},
  \texttt{school}, \texttt{user}. \\
Fehlerbehandlung &
  \texttt{GlobalExceptionHandler} mit \texttt{@RestControllerAdvice}; einheitliches
  \texttt{ErrorResponse}-DTO für alle REST-Endpunkte. \\
Authentifizierung &
  JWT-basiertes Login/Logout/Register; \texttt{JwtUtil}, \texttt{JwtAuthFilter},
  \texttt{SecurityConfig}, \texttt{AuthToken}-Persistenz. \\
School-/Klassenverwaltung &
  \texttt{School}, \texttt{SchoolClass} als JPA-Entitäten;
  eigene Controller, Service und Repository. \\
Secrets &
  JWT-Secret extern über Umgebungsvariable \texttt{\$\{JWT\_SECRET\}}; kein Plaintext im Repo. \\
\end{ssstandtabelle}

% ============================================================
\section{Chronologische Laborbuch-Einträge}
% ============================================================

Die Einträge sind in zeitlicher Reihenfolge angeordnet.
Klassifikationen: Versuchsaufbau, Messergebnis, Entscheidung, Notiz, Ergebnis, Problem.

\begin{sslabortabelle}
26.04.2026 & D0rag3on & Versuchsaufbau &
  Eintrag 01: Initialisierung des Server-Projekts unter dem vorläufigen Namen
  \texttt{SafetySpotServer}. Maven-basiertes Spring-Boot-Projekt inkl.
  \texttt{mvnw}, \texttt{pom.xml}, \texttt{persistence.xml}, \texttt{beans.xml}.
  Zu diesem Zeitpunkt noch kein funktionaler Code – lediglich ein lauffähiger,
  leerer Startpunkt als Grundlage für die Backend-Entwicklung. \\

09.06.2026 & D0rag3on & Versuchsaufbau &
  Eintrag 02: Migration des Backend-Projekts von Maven auf Gradle sowie
  Umbenennung von \texttt{SafetySpotServer} auf \texttt{ssbackend}.
  Das gesamte Projekt verwendet Gradle als einheitliches Build-System;
  Maven war ein inkonsistenter Fremdkörper.
  316 Zeilen (mvnw, pom.xml, IDE-Konfiguration) entfernt.
  Neues Modul mit Spring-Boot-Gradle-Plugin und \texttt{SsbackendApplication.java}
  aufgebaut.
  \textit{Anmerkung:} Der zugehörige erste Commit hätte früher erfolgen sollen. \\

12.06.2026 & D0rag3on & Versuchsaufbau &
  Eintrag 03: Anlegen der Paket- und Klassenstruktur für globales Error-Handling.
  Neue Pakete: \texttt{dto}, \texttt{exception}.
  Erstellt als initiale Stubs:
  \texttt{ErrorResponse} (einheitliches DTO),
  \texttt{AccessDeniedException},
  \texttt{EntityNotFoundException},
  \texttt{GlobalExceptionHandler}.
  Alle Klassen bewusst leer gehalten – Struktur vor Implementierung. \\

14.06.2026 & D0rag3on & Messergebnis &
  Eintrag 04: Vollständige Implementierung des \texttt{GlobalExceptionHandler}
  (45 Zeilen, \texttt{@RestControllerAdvice}).
  Alle definierten Ausnahmetypen werden auf HTTP-Statuscodes gemappt und als
  \texttt{ErrorResponse}-DTO zurückgegeben.
  Das globale Error-Handling ist damit abgeschlossen und einsatzbereit. \\

15.06.2026 & D0rag3on & Versuchsaufbau &
  Eintrag 05: Anlegen des vollständigen Projektskeletts für \texttt{ssbackend}:
  25 Java-Klassen über alle fachlichen Pakete verteilt.
  Pakete: \texttt{attempt} (8 Klassen), \texttt{auth} (2),
  \texttt{enums} (3), \texttt{leaderboard} (3),
  \texttt{progress} (2), \texttt{school} (3), \texttt{user} (4).
  Alle Klassen als Stubs; Ziel: Gesamtstruktur des Backends abbilden,
  bevor Implementierung beginnt. \\

22.06.2026 & D0rag3on & Ergebnis &
  Eintrag 06: JWT-Authentifizierung vollständig implementiert (PR \#31):
  \texttt{AuthController} (Login, Logout, Register), \texttt{AuthService},
  \texttt{JwtUtil}, \texttt{JwtAuthFilter}, \texttt{AuthToken} (Persistenz),
  \texttt{AuthTokenRepository}, \texttt{SecurityConfig} (stateless JWT),
  \texttt{SecurityUserDetailsService}.
  Außerdem: Foreign-Key-Beziehungen zwischen \texttt{User}, \texttt{School}
  und \texttt{SchoolClass} korrekt verknüpft. \\

22.06.2026 & D0rag3on & Notiz &
  Eintrag 07: JWT-Secret aus \texttt{application.properties} als
  Umgebungsvariable externalisiert (PR \#45).
  Der hartcodierte Wert wird jetzt aus \texttt{\$\{JWT\_SECRET\}} gelesen;
  kein sensitives Secret verbleibt im Repository.
  Entspricht dem definierten Sicherheitskonzept.
  Commit erscheint zweimal in der Historie (Cherry-Pick nach Merge-Verlust). \\

26.06.2026 & D0rag3on & Ergebnis &
  Eintrag 08: School-/Klassen-Logik vollständig implementiert (PR \#55):
  \texttt{SchoolClass}-Entität mit JPA-Beziehung zu \texttt{School};
  \texttt{SchoolClassController}, \texttt{SchoolClassService},
  \texttt{SchoolClassRepository}.
  Ergänzend: \texttt{SchoolController} und \texttt{SchoolService} fertiggestellt.
  DTOs: \texttt{CreateClass}, \texttt{CreateSchool}, \texttt{UpdateSchoolClass},
  \texttt{UpdateSchoolRequest}. \\
\end{sslabortabelle}

% ============================================================
\section{Zusammenfassung der Ergebnisse}
% ============================================================

\textbf{Fachliches Ergebnis:}
Das Backend-Grundgerüst steht: Authentifizierung, Fehlerbehandlung und
School-/Klassen-Verwaltung sind vollständig implementiert und getestet.

\textbf{Technisches Ergebnis:}
Einheitliche Gradle-Build-Struktur, JWT-basierte Stateless-Auth,
globales Error-Handling mit konsistenten HTTP-Fehlerantworten,
JPA-Entitäten für School und SchoolClass.

\textbf{Sicherheit:}
JWT-Secret wird ausschließlich über Umgebungsvariablen übergeben;
kein Plaintext-Secret im Repository.

\textbf{Lerneffekt:}
Frühe Commits erleichtern das Nachvollziehen von Änderungen erheblich.
Für zukünftige Projekte: kleinere, häufigere Commits von Beginn an.

% ============================================================
\section{Auszug aus der Git-Historie von D0rag3on}
% ============================================================

\begin{ssgittabelle}
26.04.2026 & 5b7e6c6 &
  Initialization of Server project &
  Erstes Maven-basiertes Spring-Boot-Projekt. \\
09.06.2026 & dca044a &
  Shouldve been in the last commmit &
  Restliche Dateien zur Serverinitialisierung. \\
12.06.2026 & ba0059d &
  Added new java files and packages for global Error Handling &
  \texttt{dto}/\texttt{exception}-Pakete mit Stubs. \\
14.06.2026 & e6ecf7d &
  Finished up the global error handling &
  \texttt{GlobalExceptionHandler} vollständig. \\
15.06.2026 & 368b144 &
  Added project skeleton structure, lemme cook &
  25 Java-Klassen als Stubs über alle Pakete. \\
22.06.2026 & 98c6d56 &
  Authentification and some other stuff regarding FK (\#31) &
  JWT-Auth, SecurityConfig, FK-Beziehungen. \\
22.06.2026 & 3c9dbfe &
  Added environment variable &
  JWT-Secret externalisiert. \\
26.06.2026 & dfd8c02 &
  Added Class and SchoolClass logic (\#55) &
  School, SchoolClass, Controller, Service, Repository. \\
\end{ssgittabelle}

% ============================================================
\section{Offene Punkte}
% ============================================================

\begin{sshinweistabelle}
Spring-Shell-Kommandos &
  Implementierung administrativer CLI-Kommandos über \texttt{ssshell}
  zum Anlegen und Verwalten von Demo- und Produktivdaten
  ohne direkten Datenbankzugriff. \\
Neuer Szenariotyp &
  Konzeption und Implementierung eines weiteren Szenariotyps im Backend
  (neben \texttt{HotspotItem} und \texttt{TrueFalseItem}),
  inklusive API-Endpunkte und App-Integration. \\
Test-Abdeckung &
  Unit- und Integrationstests für \texttt{SchoolClass}-API sowie
  \texttt{AuthController} ergänzen. \\
\end{sshinweistabelle}
